% Converted from the audited Markdown manuscript; responsible author: Carptopus.
\documentclass[11pt]{article}
\usepackage[a4paper,margin=28mm]{geometry}
\usepackage[T1]{fontenc}
\usepackage[utf8]{inputenc}
\usepackage{lmodern}
\usepackage{microtype}
\usepackage{amsmath,amssymb,amsthm,mathtools}
\usepackage{xcolor}
\usepackage[colorlinks=true,linkcolor=blue!55!black,citecolor=blue!55!black,urlcolor=blue!65!black]{hyperref}
\usepackage{fancyhdr}

\newtheorem{theorem}{Theorem}[section]
\newtheorem{lemma}[theorem]{Lemma}
\newtheorem{corollary}[theorem]{Corollary}
\numberwithin{equation}{section}

\pagestyle{fancy}
\fancyhf{}
\fancyhead[L]{Carptopus}
\fancyhead[R]{Odd-period quaternion ringsets}
\fancyfoot[C]{\thepage}
\setlength{\headheight}{14pt}
\setlength{\parindent}{0pt}
\setlength{\parskip}{0.55em}
\setlength{\emergencystretch}{4em}

\title{Periodic two-letter ringsets in the Lipschitz quaternions:\\
the general odd-period classification}
\author{Carptopus\\\href{mailto:carptopus@163.com}{carptopus@163.com}}
\date{24 August 2026}

\hypersetup{
  pdftitle={Periodic two-letter ringsets in the Lipschitz quaternions: the general odd-period classification},
  pdfauthor={Carptopus},
  pdfsubject={A complete ringset classification for periodic two-letter Lipschitz quaternion sets of arbitrary odd period},
  pdfkeywords={integer-valued polynomial, Lipschitz quaternion, ringset, null ideal, periodic word, fixed divisor, prime power}
}

\begin{document}
\maketitle

\begin{center}
\small\itshape
Public Beta v0.1. The theorem and the complete manuscript have passed
independent internal zero-trust audits. External mathematical review and
journal peer review are pending.
\end{center}
\vspace{0.8em}

\begin{abstract}

Let $\mathbf L=\mathbb Z[\mathbf i,\mathbf j,\mathbf k]$ be the ring of
Lipschitz quaternions. For a periodic word
$w:\mathbb Z\to\{\mathbf i,\mathbf j\}$ of an arbitrary odd period $m$, put

\[
S_w=\{a+w(a):a\in\mathbb Z\}\subseteq\mathbf L.
\]

We classify exactly when $S_w$ is a ringset, meaning that the rational
quaternion polynomials integer-valued on $S_w$ are closed under
multiplication. The criterion is that $w$ is nonconstant and, for every prime
$p\mid m$ with $p\equiv3\pmod4$, every position class modulo
$p^{v_p(m)}$ is bichromatic. This removes the squarefreeness hypothesis from
the preceding odd-period classification.

The sufficiency direction combines split prime-power cores, paired fibres at
nonsplit primes, periodic congruence orbits, and the Chinese remainder
theorem. For necessity, a monochromatic class modulo $p^\nu$ is detected in
$\mathbf L/p^{V+1}\mathbf L$, where
$V=v_p((p^\nu-1)!)$. A falling-factorial product of central quadratic
minimal polynomials gives an explicit null polynomial whose coefficientwise
right multiple is non-null. Thus the higher prime-power condition is both
necessary and sufficient; it is not inferred from finite computation.

\end{abstract}

\section{Introduction}


Let $\mathbb D$ be the rational Hamilton quaternion algebra and

\[
\mathbf L=\mathbb Z[\mathbf i,\mathbf j,\mathbf k],
\qquad
\mathbf i^2=\mathbf j^2=-1,
\qquad
\mathbf i\mathbf j=\mathbf k=-\mathbf j\mathbf i.
\]

The variable in $\mathbb D[x]$ is central, and polynomials are evaluated on
the right: if $f(x)=\sum_d c_dx^d$, then
$f(\alpha)=\sum_d c_d\alpha^d$. For $S\subseteq\mathbf L$, define

\[
\mathrm{Int}(S,\mathbf L)
=\{f\in\mathbb D[x]:f(S)\subseteq\mathbf L\}.
\]

Following Werner \cite{werner2026}, $S$ is a \emph{ringset} if
$\mathrm{Int}(S,\mathbf L)$ is a subring of $\mathbb D[x]$. In a
noncommutative coefficient ring this closure is not automatic.

Werner \cite{werner2026} classified the finite ringsets in the Lipschitz and Hurwitz
quaternion rings and established residue-null-ideal tools for infinite
sets. A preceding paper \cite{carptopus2026} applied these tools to periodic two-letter sets
and obtained a complete criterion for odd squarefree periods. The
squarefree hypothesis entered at a precise point: when $p\mid m$, it forced

\[
\gcd(p^e,m)=p
\]

for every $e\ge1$. If $p^2\mid m$, a residue fibre can instead remember a
higher position class, and a finite-field obstruction no longer detects all
failures.

The purpose of this paper is to close that boundary for every odd period.
Let $m\ge1$ be odd, let $w:\mathbb Z\to\{\mathbf i,\mathbf j\}$ satisfy
$w(a+m)=w(a)$, and put

\[
S_w=\{a+w(a):a\in\mathbb Z\}.
\]

The displayed period is not required to be least.

\begin{theorem}
\label{thm:main}

Let $m$ be odd and let $w:\mathbb Z\to\{\mathbf i,\mathbf j\}$ have
period $m$. Then $S_w$ is a ringset of $\mathbf L$ if and only if both of the
following conditions hold:

\begin{enumerate}
\item $w$ is nonconstant;
\item for every prime $p\mid m$ with $p\equiv3\pmod4$, if
$\nu=v_p(m)$, then

$$
\{w(a):a\equiv c\pmod {p^\nu}\}
=\{\mathbf i,\mathbf j\}
\qquad(c\in\mathbb Z/p^\nu\mathbb Z).
$$

\end{enumerate}

\end{theorem}

When $m$ is squarefree, $\nu=1$ and Theorem 1.1 specializes exactly to the
criterion in \cite{carptopus2026}. The new content is the necessity and sufficiency of the
highest prime-power position classes when the period is not squarefree.

\begin{corollary}
\label{cor:split-only}

If every prime divisor of the odd period $m$ is congruent to $1$ modulo $4$,
then every nonconstant binary word of period $m$ produces a ringset $S_w$.

\end{corollary}

The proof does not claim as new the residue criterion \cite{frisch2018,werner2026}, finite
matrix-ring core theory \cite{werner2022,rissnerwerner2025}, fixed-divisor methods \cite{peruginelli2014}, or the split and
paired-fibre lemmas already used in \cite{carptopus2026}. The new ingredient is an explicit
higher-prime-power obstruction tailored to a monochromatic periodic
cylinder.

\section{Null ideals and local tools}


For a ring $A$ and $T\subseteq A$, let

\[
\mathcal N(T,A)=\{f\in A[x]:f(t)=0\text{ for every }t\in T\}.
\]

This is always a left ideal. The set $T$ is \emph{core} if
$\mathcal N(T,A)$ is a two-sided ideal.

We use the following residue criterion, which follows from Werner
\cite[Lemmas 5.3 and 5.4]{werner2026}.

\begin{lemma}[Residue criterion]
\label{lem:residue}

Let $S\subseteq\mathbf L$, and let $\overline S_n$ be its image in
$\mathbf L/n\mathbf L$.

\begin{enumerate}
\item If $\overline S_n$ is core for every $n\ge2$, then $S$ is a ringset.
\item If, for some $n$, there are
$g\in\mathcal N(\overline S_n,\mathbf L/n\mathbf L)$ and
$u\in\{\mathbf i,\mathbf j,\mathbf k\}$ such that
$gu\notin\mathcal N(\overline S_n,\mathbf L/n\mathbf L)$, then $S$
is not a ringset.

\end{enumerate}

\end{lemma}

Here $gu$ means uniform right multiplication of all coefficients by $u$.

We also use two local core lemmas from \cite{carptopus2026}. They are restated to make the
argument self-contained.

\begin{lemma}[Paired fibres]
\label{lem:paired}

Let $q\ge2$. Any subset

\[
T\subseteq
\{a+\mathbf i,a+\mathbf j:a\in\mathbb Z/q\mathbb Z\}
\]

that contains both letters above every real residue is core in
$\mathbf L/q\mathbf L$.

\begin{proof}

For a fixed real residue $a$, divide a null polynomial $f$ by the central
minimal polynomial

\[
\mu_a(x)=x^2-2ax+a^2+1.
\]

The remainder is $r(x)=\gamma_1x+\gamma_0$. Its vanishing at
$a+\mathbf i$ and $a+\mathbf j$ gives
$\gamma_1(\mathbf i-\mathbf j)=0$, hence $2\gamma_1=0$ because
$(\mathbf i-\mathbf j)(-\mathbf i+\mathbf j)=2$. For any quaternion basis
unit $u$ and either root $\alpha$,

\[
(ru)(\alpha)-r(\alpha)u
=\gamma_1(u\alpha-\alpha u)=0,
\]

since the commutator is divisible by $2$. Centrality of $\mu_a$ passes this
from the remainder to $f$. Thus the null ideal is stable under right
multiplication by constants, and therefore by all polynomials because it is
a left ideal and $x$ is central.
\end{proof}

\end{lemma}


\begin{lemma}[Split one-letter fibres]
\label{lem:split}

Let $p\equiv1\pmod4$, $e\ge1$, and
$R=\mathbb Z/p^e\mathbb Z$. If

\[
T\subseteq\{a+\mathbf i,a+\mathbf j:a\in R\}
\]

contains at least one point above every real residue, then $T$ is core in
$\mathbf L/p^e\mathbf L$.

\begin{proof}

Choose $s\in R^\times$ with $s^2=-1$. The standard representation

\[
\mathbf i\longmapsto
\begin{pmatrix}s&0\\0&-s\end{pmatrix},
\qquad
\mathbf j\longmapsto
\begin{pmatrix}0&1\\-1&0\end{pmatrix}
\]

identifies $\mathbf L/p^e\mathbf L$ with $M_2(R)$. For
$u\in\{\mathbf i,\mathbf j\}$, put

\[
e_u^+=\frac{1+u/s}{2},
\qquad
e_u^-=\frac{1-u/s}{2}.
\]

These are complementary orthogonal idempotents and

\[
(a+u)^d=(a+s)^d e_u^+ +(a-s)^d e_u^-.
\]

The image of a plus idempotent for either letter and the image of a minus
idempotent for either letter are complementary rank-one summands. Explicit
generators give determinants $1,-s,1,-2s$ in the four possible pairings,
all units in $R$.

For a null polynomial $f(x)=\sum_dc_dx^d$, let
$F(z)=\sum_dc_dz^d$. At a fixed $z\in R$, use a selected point of real part
$z-s$ and one of real part $z+s$. Their zero evaluations give
$F(z)e_{u_1}^+=F(z)e_{u_2}^-=0$, so $F(z)=0$. Consequently the scalar
polynomial for every coefficientwise right multiple $fr$ is
$F(z)r=0$, and the idempotent decomposition shows that $fr$ is null on
$T$. The null ideal is two-sided.
\end{proof}

\end{lemma}


\section{Periodic congruence fibres}


\begin{lemma}[Periodic orbit lemma]
\label{lem:orbit}

Let $w$ have period $m$, let $q\ge1$, and fix
$a\in\mathbb Z/q\mathbb Z$. The positions modulo $m$ attained by integers
congruent to $a$ modulo $q$ are

\[
a+q\mathbb Z\pmod m
=\{r\in\mathbb Z/m\mathbb Z:r\equiv a
  \pmod{\gcd(q,m)}\}.
\]

\begin{proof}

The subgroup generated by $q$ in $\mathbb Z/m\mathbb Z$ is the subgroup of
multiples of $\gcd(q,m)$. Translate it by $a$.
\end{proof}

\end{lemma}


For $q=p^e$ and $\nu=v_p(m)$, a real residue fibre therefore sees one
position class modulo

\[
p^{\min(e,\nu)}.
\]

If every class modulo $p^\nu$ is bichromatic, every coarser class is a union
of bichromatic classes and is also bichromatic.

The criterion is independent of using a nonminimal period. Indeed, if the
least period is $d\mid m$, then a class modulo $p^{v_p(m)}$ projects onto a
class modulo $p^{v_p(d)}$; its set of letters is exactly the latter class's
set of letters.

\section{A prime-power fixed-divisor lemma}


The necessity argument requires a precise valuation dichotomy.

\begin{lemma}
\label{lem:fixed-divisor}

Let $p$ be a prime, $\nu\ge1$, and put

\[
N=p^\nu-1,
\qquad
V=v_p(N!),
\qquad
P(d)=\prod_{r=1}^{N}(d-r).
\]

With the convention $v_p(0)=+\infty$,

\[
v_p(P(d))=
\begin{cases}
V,&p^\nu\mid d,\\
\ge V+1,&p^\nu\nmid d.
\end{cases}
\]

\begin{proof}

If $p^\nu\mid d$, then $v_p(d-r)=v_p(r)$ for every
$1\le r<p^\nu$, and summing gives $V$.

Otherwise choose $s\in\{1,\ldots,N\}$ with
$s\equiv d\pmod {p^\nu}$. The factor $d-s$ has valuation at least $\nu$.
The nonzero residue classes of the remaining $d-r$ modulo $p^\nu$ are the
nonzero classes other than $s$, up to permutation. Hence

\[
v_p(P(d))\ge V-v_p(s)+\nu\ge V+1,
\]

because $v_p(s)\le\nu-1$.
\end{proof}

\end{lemma}


We shall work modulo $p^E$, where

\[
E=V+1.
\]

Notice that $E\ge\nu$: the assertion is immediate for $\nu=1$, while for
$\nu\ge2$ the factors $p,p^2,\ldots,p^{\nu-1}$ already contribute at least
$\nu-1$ to $V$.

\section{A higher nonsplit obstruction}


Let $p\equiv3\pmod4$ and $\nu\ge1$. Suppose that a position class
$c\pmod {p^\nu}$ is monochromatic. We construct a finite quotient witness
without assuming any classification of subsets over $\mathbb Z/p^e\mathbb Z$.

Assume first that the class contains only $\mathbf i$; the other case is
symmetric. Use $N,V,E$ from Section 4 and define, in
$(\mathbf L/p^E\mathbf L)[x]$,

\[
\mu_{c+r}(x)=(x-(c+r))^2+1
\qquad(1\le r\le N),
\]

\[
Q_c(x)=\prod_{r=1}^{N}\mu_{c+r}(x),
\qquad
h_c(x)=(x-(c+\mathbf i))Q_c(x).
\]

Every $\mu_{c+r}$, and hence $Q_c$, has central coefficients.

For $q\in\mathbf L$, write

\[
v_p(q)=\max\{t\ge0:q\in p^t\mathbf L\},
\]

with $v_p(0)=+\infty$. Equivalently, this is the minimum of the $p$-adic
valuations of the four Lipschitz coordinates. Multiplication by a quaternion
whose norm is prime to $p$ preserves this valuation.

\begin{lemma}
\label{lem:obstruction}

The polynomial $h_c$ is null on the image of $S_w$ modulo $p^E$, but the
coefficientwise right multiple $h_c\mathbf j$ is not null.

\begin{proof}

Since $E\ge\nu$ and $\nu=v_p(m)$, we have
$\gcd(p^E,m)=p^\nu$. Thus Lemma 3.1 says that every real fibre modulo $p^E$
sees exactly one position class modulo $p^\nu$.

Take a point $y=a+u$ of the periodic image, where
$u\in\{\mathbf i,\mathbf j\}$, and write $d=a-c$. Direct calculation gives

\[
\mu_{c+r}(y)=(d-r)(d-r+2u).
\]

The second factor is a unit modulo $p^E$. Indeed, its norm modulo $p$ is
$(d-r)^2+4$, which cannot vanish because $-1$ is not a square when
$p\equiv3\pmod4$. Therefore

\[
v_p(Q_c(y))=v_p(P(d)).
\]

If $p^\nu\nmid d$, Lemma 4.1 shows that $Q_c(y)=0$ modulo $p^E$. If
$p^\nu\mid d$ and $u=\mathbf i$, the first factor of $h_c(y)$ is $d$ and
the total valuation is at least $V+\nu\ge E$. If
$p^\nu\mid d$ and $u=\mathbf j$, then the periodic orbit lemma says that
the relevant position class is the monochromatic class $c$, so no such
point occurs. Thus $h_c$ is null on the entire periodic image.

Because $Q_c$ is central, coefficientwise right multiplication gives

\[
h_c\mathbf j=((x-(c+\mathbf i))\mathbf j)Q_c.
\]

At the point $c+\mathbf i$ the first factor evaluates, with right
evaluation, to

\[
\mathbf j(c+\mathbf i)-(c+\mathbf i)\mathbf j=-2\mathbf k,
\]

a $p$-adic unit. Lemma 4.1 gives
$v_p(Q_c(c+\mathbf i))=V$. Hence

\[
(h_c\mathbf j)(c+\mathbf i)\ne0\pmod {p^{V+1}}.
\]

This proves the claim.
\end{proof}

\end{lemma}


If the monochromatic letter is $\mathbf j$, replace the leading factor by
$x-(c+\mathbf j)$ and right-multiply coefficients by $\mathbf i$.

\section{Proof of Theorem 1.1}

\begin{proof}

\medskip
\noindent\textbf{Sufficiency.}

Assume conditions 1 and 2. It is enough to show that the image of $S_w$ is
core modulo every prime power $p^e$.

If $p\equiv1\pmod4$, every real residue occurs and Lemma 2.3 applies; no
bichromatic condition is needed.

If $p=2$, then $\gcd(2^e,m)=1$ because $m$ is odd. Every real fibre traverses
the whole period, so nonconstancy makes it bichromatic. Lemma 2.2 applies.

Now let $p\equiv3\pmod4$. If $p\nmid m$, the same coprimality argument and
Lemma 2.2 apply. If $\nu=v_p(m)>0$, Lemma 3.1 says that a real fibre modulo
$p^e$ sees one position class modulo $p^{\min(e,\nu)}$. Condition 2 makes
every class at level $p^\nu$ bichromatic, and therefore every coarser class
is bichromatic. Lemma 2.2 again applies.

For an arbitrary denominator $n=\prod_rp_r^{e_r}$, the Chinese remainder
theorem gives

\[
\mathbf L/n\mathbf L
\simeq\prod_r\mathbf L/p_r^{e_r}\mathbf L.
\]

Polynomial coefficients and right evaluation are componentwise, so the null
ideal of an image is the product of the null ideals of its projections. Each
prime-power projection is core; hence the image modulo $n$ is core. Lemma
2.1 proves that $S_w$ is a ringset.

\medskip
\noindent\textbf{Necessity.}

If $w$ is constant, reduce modulo $3$. The finite-field missing-letter
obstruction from \cite[Lemma 5.1]{carptopus2026} supplies a null polynomial whose right
multiple is non-null. Lemma 2.1 shows that $S_w$ is not a ringset.

Suppose next that condition 2 fails. Choose $p\mid m$ with
$p\equiv3\pmod4$, put $\nu=v_p(m)$, and choose a monochromatic class
$c\pmod {p^\nu}$. Lemma 5.1 gives a null polynomial in the image modulo
$p^{V+1}$ whose coefficientwise right multiple is non-null. The second part
of Lemma 2.1 again shows that $S_w$ is not a ringset. This proves both
necessary conditions and completes the theorem.
\end{proof}


\section{Consequences and examples}


Theorem 1.1 removes the entire odd-period boundary left open in \cite{carptopus2026}. It also
shows exactly why checking classes only modulo $p$ is insufficient when
$p^2\mid m$.

For a displayed period $m=9$, it is possible for all three classes modulo
$3$ to be bichromatic. Nevertheless every class modulo $9$ consists of a
single word position and is therefore monochromatic. Thus no period-$9$
word satisfies the theorem: choosing any one of these highest classes gives
an obstruction modulo $27$, since

\[
v_3(8!)=2,
\qquad E=3.
\]

For period $m=45=3^2\cdot5$, only the classes modulo $9$ impose a colour
condition; the split prime $5$ imposes none. For a period divisible only by
primes congruent to $1$ modulo $4$, Corollary 1.2 gives all nonconstant
words, regardless of repeated prime factors.

The criterion can be checked using any stated odd period. If a larger
nonminimal period is used, the extra prime-power classes merely repeat the
classes belonging to the least period, as observed after Lemma 3.1.

\section{Computational calibration}


The proof is symbolic. Computation was used only to discover and test the
prime-power boundary.

\begin{itemize}
\item For one coarse class modulo $p$ in the quotient modulo $p^2$, each of its
$p$ lifts was allowed to contain only $\mathbf i$, only $\mathbf j$, or
both. For $p=3$, all $3^3=27$ patterns agreed with the local lift-count
rule that each letter must occur in at least two lifts. This is weaker than
the hypothesis of Lemma 2.2 and is only a finite local-ring observation.
\item For $p=7$ modulo $49$, samples agreed with the same lift-count threshold:
deleting $\mathbf j$ from one or five lifts remained stable, while deleting
it from six or seven lifts did not.
\item Explicit obstruction polynomials were evaluated for
$(p,\nu)=(3,1),(3,2),(7,1)$.
\item An independent row-module computation over $\mathbb Z/27\mathbb Z$
rejected the finite fibre pattern having only $\mathbf i$ above real
residues congruent to $0$ modulo $9$ and both letters elsewhere. This
finite quotient pattern is not itself the residue image of a two-letter
word of displayed period $9$.

\end{itemize}

These checks calibrate the formulas but do not prove the general theorem.
The arbitrary-$p,\nu$ conclusion comes from Lemmas 4.1 and 5.1.

\section{Prior-work and novelty boundary}


The general ringset/null-ideal link is due to Frisch \cite{frisch2018}. Werner \cite{werner2022}
classified null ideals of subsets of matrix rings over fields, and Rissner
and Werner \cite{rissnerwerner2025} studied their enumeration. Werner \cite{werner2026} classified finite
Lipschitz and Hurwitz quaternion ringsets and established the residue and
right-multiplication criteria used here. Peruginelli \cite{peruginelli2014} developed
prime-power fixed-divisor ideals in the commutative integer-valued-polynomial
setting. None of these general tools is claimed as new.

The preceding paper \cite{carptopus2026} proves the squarefree instance of Theorem 1.1 and
contains the split and finite-field local mechanisms reused here. The present
paper claims only the removal of squarefreeness, the exact
$p^{v_p(m)}$ condition, and the explicit higher-prime-power obstruction.
The prime-square analysis is an intermediate part of this general result,
not a separate result of comparable scope.

As of the bounded literature search completed on 23 August 2026, no public
source was found that directly states this all-odd-period classification.
That assessment is not a certificate of absolute global priority. The result
has not received external expert review or formal peer review.

\section*{Declaration of generative AI and AI-assisted technologies}


During the research and preparation of this work, the author used OpenAI
Codex to assist with literature searches, symbolic exploration, exploratory
code review, proof checking, and language drafting. The author reviewed and
edited all generated output and takes full responsibility for the content of
this manuscript.

\begin{thebibliography}{9}

\bibitem{carptopus2026}
Carptopus,
\emph{Periodic infinite ringsets in the Lipschitz quaternions},
Public Beta v0.1-beta (2026).
\href{https://doi.org/10.5281/zenodo.22069614}{doi:10.5281/zenodo.22069614}.

\bibitem{peruginelli2014}
G.~Peruginelli,
\emph{Primary decomposition of the ideal of polynomials whose fixed divisor is divisible by a prime power},
J. Algebra 398 (2014), 227--242.
\href{https://doi.org/10.1016/j.jalgebra.2013.09.016}{doi:10.1016/j.jalgebra.2013.09.016};
\href{https://arxiv.org/abs/1304.7450}{arXiv:1304.7450}.

\bibitem{frisch2018}
S.~Frisch,
\emph{Polynomial functions on subsets of non-commutative rings---a link between ringsets and null-ideal sets},
ITM Web Conf. 20 (2018), 01003.
\href{https://doi.org/10.1051/itmconf/20182001003}{doi:10.1051/itmconf/20182001003}.

\bibitem{werner2022}
N.~J. Werner,
\emph{Null ideals of subsets of matrix rings over fields},
Linear Algebra Appl. 642 (2022), 50--72.
\href{https://doi.org/10.1016/j.laa.2022.02.007}{doi:10.1016/j.laa.2022.02.007}.

\bibitem{werner2026}
N.~J. Werner,
\emph{Integer-valued polynomials on subsets of quaternion algebras},
J. Algebra 686 (2026), 195--219.
\href{https://doi.org/10.1016/j.jalgebra.2025.08.011}{doi:10.1016/j.jalgebra.2025.08.011};
\href{https://arxiv.org/abs/2412.20609}{arXiv:2412.20609v2}.

\bibitem{rissnerwerner2025}
R.~Rissner and N.~J. Werner,
\emph{Counting core sets in matrix rings over finite fields},
Linear Algebra Appl. 720 (2025), 1--25.
\href{https://doi.org/10.1016/j.laa.2025.04.006}{doi:10.1016/j.laa.2025.04.006};
\href{https://arxiv.org/abs/2405.04106}{arXiv:2405.04106v2}.

\end{thebibliography}

\end{document}
